\scanpage{66}
\[
\handwritten{\xi\in}H^{\handwritten{2}}(k_{\mathrm{pl}},G)
\]
est la classe de la gerbe tannakienne qui définit $\handwritten{\unsure{\mathcal{M}|\deleted{M}\mathcal{M}|\deleted{\mathcal{A}}\mathcal{M}}}$. Donc la connaissance du caractère numérique de $\handwritten{\mathcal{M}}$ équivaut à la donnée, en plus d\unsure{\deleted{e}u|\deleted{u}e}\\ \hspace*{3cm}\deleted{\inserted{(qui implique déjà la donnée}}\\ groupe de t.m. $G$ sur $k$, des classes $u_{\unsure{\deleted{\unsure{\xi|\sigma}}\handwritten{*}|\handwritten{\xi}|\handwritten{\sigma}}}(\handwritten{\xi})$ pour tous les homomorphismes $u$ de $G$ dans des tores de la forme $\handwritten{\prod}_{Z/k}\underline{G}_m$, où $Z$ est une extension finie de $k$. \deleted{(}On vient de voir en effet que la donnée du caractère numérique permet de reconstituer ce qui précède. Pour voir l'inverse, notons que \deleted{si} le rang $\unsure{\deleted{d}n|\deleted{n}d}(\handwritten{\deleted{\unsure{\xi}}\sigma})$ de $M_{\handwritten{\sigma}}$ est déterminé par la formule
\[
n(\handwritten{\sigma})=\operatorname{card}(\handwritten{\sigma})\operatorname{\unsure{\deleted{d}o|\deleted{o}d}rd}(\handwritten{\xi_\sigma})
\]
(où $\operatorname{card}\handwritten{\sigma}$ est le cardinal de l'orbite $\handwritten{\sigma}$, $\operatorname{ord}(\handwritten{\xi_\sigma})$ est l'ordre de $\handwritten{\xi_\sigma}$ dans $\operatorname{Br}(Z_{\handwritten{\sigma}})$).

La question de savoir si le caractère numérique détermine à équivalence de catégories tannakiennes près revient donc, pour $G$ fixé, à la question de savoir si un élément \deleted{de} $\handwritten{\xi\in}H^{\handwritten{2}}(k_{\mathrm{pl}},G_{\deleted{\times}})$ est connu quand on connait les $u_{\handwritten{\unsure{*|\sigma}}}(\handwritten{\xi})$ pour $u:G\handwritten{\unsure{\longrightarrow|\rightharpoonup}}\underline{Z}^{\unsure{\times|x|\mathbb{x}|*}}$ comme dessus. \deleted{On va ete\unsure{s}} C'est le cas pour $G$ diagonalisable \deleted{pro-dénombrable} (pro-dénombrable ?). On va examiner deux autres cas.

\uline{Exemple} \unsure{\deleted{G}1|\deleted{1}G}. $G$ est un tore de dimension 1, donc (s'il n'est trivial) tordu par une extension quadratique $Z$ de $k$. On a donc une suite exacte
\diagram{\begin{tikzcd}[ampersand replacement=\&]
0 \arrow[r] \& G \arrow[r,"u"] \& \underline{Z}^{\times} \arrow[r,"{\unsure{\mathbb{N}|N}_{Z/k}}"] \& \underline{G}_m \arrow[r] \& 0\ ,
\end{tikzcd}}
d'où on conclut, par la suite exacte de cohomologie et le th. 90 que $\operatorname{Ker}(u:H^{\handwritten{2}}(k,G)\handwritten{\longrightarrow}H^{\handwritten{2}}(\unsure{\deleted{i}k|\deleted{k}i},\underline{Z}^{\times}))$ est nul. Donc dans le cas envisagé, le caractère numérique détermine $\handwritten{\mathcal{M}}$ à équivalence près.

\deleted{\uline{Exemple} 2. Gardons \deleted{l\unsure{'|i|t}e} $G$ comme dessus, soit $n$ un entier $\quad2$, et soit $G'=\operatorname{Ker}n.\operatorname{id}_G$. La suite exacte}
\diagram{\begin{tikzcd}[ampersand replacement=\&]
\deleted{0} \& \deleted{G'} \& \deleted{G} \arrow[r,phantom,"\deleted{n}"] \& \deleted{G} \& \deleted{0}
\end{tikzcd}}
\deleted{et le résultat précé\unsure{\deleted{e}d|\deleted{d}e}ent montre que l'ensemble des classes de}

% archive sheet number: 65
